\section{Maximum XOR With an Element from Array (Queries)}
\label{sec:trie-max-xor-query}

\problemstatement{Given an array $\textit{nums}$ and a list of queries,
where each query is a pair $(x_i, m_i)$, answer for each query the
maximum value of $x_i \oplus \textit{nums}[j]$ over all $j$ with
$\textit{nums}[j] \le m_i$. If no element of $\textit{nums}$ satisfies
$\textit{nums}[j] \le m_i$, the answer for that query is $-1$.}

\begin{patternbox}
This is Section~\ref{sec:trie-max-xor}'s binary trie with one new
constraint layered on top: each query only considers a \textbf{subset}
of $\textit{nums}$ (those $\le m_i$), and the subset differs per query.
The keyword combination \emph{``maximum XOR'' plus ``subject to a
threshold that varies per query''} signals an \textbf{offline} technique:
sort the queries by their threshold and process them in increasing
order, incrementally growing the trie to include exactly the elements
eligible for each query, so that no element is ever inserted before it
is allowed to be, and no element needs to be removed once inserted.
\end{patternbox}

\subsection*{Understanding the problem carefully}

If queries arrived with arbitrary, unsorted thresholds $m_i$, we would
need to repeatedly figure out ``which elements of $\textit{nums}$ are
$\le m_i$'' for each query independently. But if we instead
\textbf{sort the queries by} $m_i$ and process them in increasing order,
the set of eligible elements only ever \textbf{grows} from one query to
the next (a higher threshold can only admit more elements, never fewer)
-- so we can maintain a single pointer into a separately sorted copy of
$\textit{nums}$, inserting new elements into the trie as the threshold
grows, and never needing to remove anything.

\begin{ideabox}[Offline: Sort Both Queries and the Array by Threshold]
Sort $\textit{nums}$ in increasing order. Sort the queries by $m_i$,
remembering each query's original index (to place answers back
correctly). Maintain a pointer $j$ into the sorted $\textit{nums}$,
starting at $0$. For each query, in increasing order of $m_i$: while $j$
has not exhausted $\textit{nums}$ and $\textit{nums}[j] \le m_i$, insert
$\textit{nums}[j]$ into the binary trie and advance $j$. Then, if the
trie is non-empty, answer this query using
Section~\ref{sec:trie-max-xor}'s $\textit{maxXorWith}(x_i)$; otherwise
(no eligible elements at all), answer $-1$.
\end{ideabox}

\subsection*{Why processing queries in sorted threshold order is correct and never wastes work}

Because eligibility only grows as the threshold increases, every element
inserted into the trie while answering an earlier (smaller-threshold)
query remains correctly eligible for every later (larger-threshold)
query too -- nothing inserted ever needs to be removed. This means the
pointer $j$ only ever moves forward, each element of $\textit{nums}$ is
inserted into the trie exactly once across the entire batch of queries,
and the total insertion work across all queries combined is $O(n)$
insertions, not $O(nq)$ for $q$ queries -- turning what looks like a
per-query rebuild into a single shared incremental structure.

\subsection*{Algorithm}

\begin{enumerate}
  \item Sort $\textit{nums}$ in increasing order.
  \item Sort the queries by $m_i$, keeping track of each query's
        original position.
  \item Initialize an empty binary trie and a pointer $j \gets 0$.
  \item For each query $(x_i, m_i)$, in sorted order:
  \begin{enumerate}
    \item While $j<n$ and $\textit{nums}[j] \le m_i$: insert
          $\textit{nums}[j]$ into the trie; $j \gets j+1$.
    \item If the trie is empty: this query's answer is $-1$.
    \item Else: this query's answer is
          $\textit{maxXorWith}(x_i)$.
  \end{enumerate}
  \item Place each computed answer back at its query's original
        position; return the answers.
\end{enumerate}

\subsection*{Dry run}

\dryrun{Take $\textit{nums} = [0,1,2,3,4]$ and queries
$[(3,1), (1,3), (5,6)]$.}

Sorted $\textit{nums}$ unchanged: $[0,1,2,3,4]$. Sorted queries by
$m_i$: $(3,1), (1,3), (5,6)$ (already in this order).

\begin{itemize}
  \item Query $(x{=}3, m{=}1)$: insert $\textit{nums}[j]$ while
        $\le1$: insert $0$ ($j\to1$), insert $1$ ($j\to2$); next is
        $2>1$, stop. Trie has $\{0,1\}$.
        $\textit{maxXorWith}(3)$: $3\oplus0=3$, $3\oplus1=2$; best is
        $3$.
  \item Query $(x{=}1, m{=}3)$: insert while $\le3$: insert $2$
        ($j\to3$), insert $3$ ($j\to4$); next is $4>3$, stop. Trie has
        $\{0,1,2,3\}$. $\textit{maxXorWith}(1)$: $1\oplus0=1$,
        $1\oplus1=0$, $1\oplus2=3$, $1\oplus3=2$; best is $3$.
  \item Query $(x{=}5, m{=}6)$: insert while $\le6$: insert $4$
        ($j\to5$); $j=n$, stop. Trie has $\{0,1,2,3,4\}$.
        $\textit{maxXorWith}(5)$: best pairing is $5\oplus2=7$ (checking
        all: $5{=}101$, $0{=}000\to101{=}5$, $1{=}001\to100{=}4$,
        $2{=}010\to111{=}7$, $3{=}011\to110{=}6$, $4{=}100\to001{=}1$);
        best is $7$.
\end{itemize}

Answers in original query order: $[3,3,7]$.

\subsection*{C++ implementation}

\begin{lstlisting}
#include <bits/stdc++.h>
using namespace std;

struct TrieNode {
    TrieNode* children[2] = {nullptr, nullptr};
};

class BinaryTrie {
    TrieNode* root = new TrieNode();
    static const int BITS = 32;
    bool empty = true;

public:
    void insert(int num) {
        empty = false;
        TrieNode* node = root;
        for (int b = BITS - 1; b >= 0; b--) {
            int bit = (num >> b) & 1;
            if (!node->children[bit]) node->children[bit] = new TrieNode();
            node = node->children[bit];
        }
    }

    bool isEmpty() { return empty; }

    int maxXorWith(int num) {
        TrieNode* node = root;
        int result = 0;
        for (int b = BITS - 1; b >= 0; b--) {
            int bit = (num >> b) & 1;
            int desired = 1 - bit;
            if (node->children[desired]) {
                result |= (1 << b);
                node = node->children[desired];
            } else {
                node = node->children[bit];
            }
        }
        return result;
    }
};

vector<int> maximizeXor(vector<int>& nums, vector<vector<int>>& queries) {
    sort(nums.begin(), nums.end());

    int q = queries.size();
    vector<int> order(q);
    for (int i = 0; i < q; i++) order[i] = i;
    sort(order.begin(), order.end(), [&](int a, int b) {
        return queries[a][1] < queries[b][1];
    });

    BinaryTrie trie;
    vector<int> answers(q);
    int j = 0, n = nums.size();

    for (int idx : order) {
        int x = queries[idx][0], m = queries[idx][1];
        while (j < n && nums[j] <= m) {
            trie.insert(nums[j]);
            j++;
        }
        answers[idx] = trie.isEmpty() ? -1 : trie.maxXorWith(x);
    }
    return answers;
}

int main() {
    vector<int> nums = {0, 1, 2, 3, 4};
    vector<vector<int>> queries = {{3,1}, {1,3}, {5,6}};

    for (int ans : maximizeXor(nums, queries)) cout << ans << " ";
    cout << endl;
    return 0;
}
\end{lstlisting}

\begin{complexitybox}
\textbf{Time Complexity.} Sorting $\textit{nums}$ and the queries costs
$O(n \log n + q \log q)$. Each element of $\textit{nums}$ is inserted
into the trie exactly once across the whole batch, costing
$O(n \cdot \textit{BITS})$ total. Each query performs one
$O(\textit{BITS})$ trie walk. Overall:
\[
  O((n+q)\log(n+q) + (n+q)\cdot \textit{BITS}).
\]
\textbf{Space Complexity.} $O(n \cdot \textit{BITS})$ for the trie, plus
$O(q)$ for the answers and sort order.
\end{complexitybox}

\begin{takeawaybox}
This closing problem of the chapter -- and of this entire multi-chapter
project -- is a deliberate capstone combining two ideas: the binary trie
from Section~\ref{sec:trie-max-xor}, and the classic \textbf{offline
query} technique of sorting queries by their constraint and processing
them alongside a sorted version of the data, so that a single
incrementally-built structure serves every query without ever needing to
be rebuilt or have elements removed. Whenever queries each restrict
attention to a growing prefix of sorted data, check whether sorting the
queries themselves by that same threshold turns many independent queries
into one shared incremental pass.
\end{takeawaybox}
